\section{Análisis de las funciones}

\subsection{Definición y visualización}

Las tres funciones se evaluaron sobre una malla regular de $401\times401$
puntos en $[-10,10]^2$, construida con \texttt{torch.meshgrid}. Para cada
una se generaron curvas de nivel (\texttt{contourf}) y una superficie 3D
mediante la función auxiliar \texttt{plot\_function\_2d\_3d} del notebook.

\begin{align}
f_0(x,y) &= x^2+y^2, \\
f_1(x,y) &= -20\exp\!\left[-0.2\sqrt{0.5(x^2+y^2)}\right]
  - \exp\!\left[0.5(\cos 2\pi x+\cos2\pi y)\right] + e + 20, \\
f_2(x,y) &= (x^2+y-11)^2+(x+y^2-7)^2.
\end{align}

% TODO: Insertar aquí las figuras exportadas a results/figures/ (curvas de
% nivel y superficie 3D) para f0, f1 y f2.

\subsection{Convexidad}

$f_0$ es una función cuadrática con Hessiana constante $2I$, positiva
definida en todo el dominio, por lo que es convexa. $f_1$ (Ackley) y $f_2$
(Himmelblau) no son convexas: ambas presentan múltiples mínimos locales y
puntos silla, según se muestra a continuación.

\subsection{Derivación analítica de los gradientes}
\label{sec:gradientes_analiticos}

La búsqueda de puntos críticos se apoya en \texttt{torch.autograd}, que
calcula el gradiente de forma automática. Como verificación independiente
de esos valores, a continuación se deriva cada gradiente de forma cerrada
mediante la regla de la cadena.

\subsubsection{$f_0(x,y)=x^2+y^2$}
Al ser una suma de potencias, sus derivadas parciales son inmediatas:
\[
\frac{\partial f_0}{\partial x} = 2x, \qquad \frac{\partial f_0}{\partial y} = 2y.
\]

\subsubsection{$f_1(x,y)$ (Ackley)}

Conviene separar $f_1$ en dos términos que dependen de $(x,y)$ y una
constante:
\[
f_1 = \underbrace{-20\exp\!\left[-0.2\sqrt{0.5(x^2+y^2)}\right]}_{A(x,y)}
\;\underbrace{-\exp\!\left[0.5(\cos 2\pi x+\cos2\pi y)\right]}_{B(x,y)}
\;+\;\underbrace{e+20}_{\text{constante}}.
\]
La constante no depende de $x,y$, por lo que su derivada es $0$.

\paragraph{Término $A$.} Sea $u(x,y)=0.5(x^2+y^2)$, de modo que
$A=-20\exp[-0.2\sqrt{u}]$. Aplicando la regla de la cadena de adentro
hacia afuera respecto a $x$:
\[
\frac{\partial u}{\partial x} = x,
\qquad
\frac{\partial \sqrt{u}}{\partial x} = \frac{1}{2\sqrt{u}}\cdot\frac{\partial u}{\partial x} = \frac{x}{2\sqrt{u}},
\]
\[
\frac{\partial}{\partial x}\big(-0.2\sqrt{u}\big) = \frac{-0.1\,x}{\sqrt{u}},
\qquad
\frac{\partial}{\partial x}\exp\!\left[-0.2\sqrt{u}\right] = \exp\!\left[-0.2\sqrt{u}\right]\cdot\frac{-0.1\,x}{\sqrt{u}}.
\]
Multiplicando por el $-20$ que multiplica a todo el término, los dos
signos negativos se cancelan:
\[
\frac{\partial A}{\partial x} = -20\cdot\frac{-0.1\,x}{\sqrt{u}}\exp\!\left[-0.2\sqrt{u}\right]
= \frac{2x}{\sqrt{u}}\exp\!\left[-0.2\sqrt{u}\right].
\]
Por la simetría entre $x$ y $y$ dentro de $u$, el mismo procedimiento da
\[
\frac{\partial A}{\partial y} = \frac{2y}{\sqrt{u}}\exp\!\left[-0.2\sqrt{u}\right].
\]

\paragraph{Término $B$.} Sea $v(x,y)=0.5(\cos 2\pi x+\cos 2\pi y)$, de
modo que $B=-\exp[v]$. Usando $\frac{d}{dx}\cos(2\pi x)=-2\pi\sin(2\pi
x)$ y que $\cos(2\pi y)$ no depende de $x$:
\[
\frac{\partial v}{\partial x} = 0.5\cdot(-2\pi\sin 2\pi x) = -\pi\sin(2\pi x),
\qquad
\frac{\partial}{\partial x}\exp[v] = \exp[v]\cdot(-\pi\sin 2\pi x).
\]
Multiplicando por el $-1$ que multiplica a todo el término, los signos
vuelven a cancelarse:
\[
\frac{\partial B}{\partial x} = \pi\sin(2\pi x)\exp[v].
\]
De forma análoga,
\[
\frac{\partial B}{\partial y} = \pi\sin(2\pi y)\exp[v].
\]

\paragraph{Gradiente completo.} Sumando $A$ y $B$, y recordando que la
constante $e+20$ no aporta:
\begin{align}
\frac{\partial f_1}{\partial x} &= \frac{2x}{\sqrt{0.5(x^2+y^2)}}\exp\!\left[-0.2\sqrt{0.5(x^2+y^2)}\right]
  + \pi\sin(2\pi x)\exp\!\left[0.5(\cos 2\pi x+\cos2\pi y)\right], \\
\frac{\partial f_1}{\partial y} &= \frac{2y}{\sqrt{0.5(x^2+y^2)}}\exp\!\left[-0.2\sqrt{0.5(x^2+y^2)}\right]
  + \pi\sin(2\pi y)\exp\!\left[0.5(\cos 2\pi x+\cos2\pi y)\right].
\end{align}
Estas expresiones no están definidas en $(0,0)$ (división entre
$\sqrt{u}=0$), lo cual coincide con la falta de suavidad del mínimo
global de Ackley señalada más adelante. Se verificó de forma
independiente con WolframAlpha que ambas derivadas coinciden con esta
forma cerrada.

\subsubsection{$f_2(x,y)$ (Himmelblau)}

$f_2$ es la suma de dos términos al cuadrado; cada uno se deriva con la
regla de la cadena $\frac{d}{dx}[(\cdot)^2]=2(\cdot)\cdot(\cdot)'$.

Respecto a $x$:
\[
\frac{\partial}{\partial x}\big[(x^2+y-11)^2\big] = 2(x^2+y-11)\cdot 2x = 4x(x^2+y-11),
\]
\[
\frac{\partial}{\partial x}\big[(x+y^2-7)^2\big] = 2(x+y^2-7)\cdot 1 = 2(x+y^2-7),
\]
\[
\frac{\partial f_2}{\partial x} = 4x(x^2+y-11) + 2(x+y^2-7).
\]

Respecto a $y$:
\[
\frac{\partial}{\partial y}\big[(x^2+y-11)^2\big] = 2(x^2+y-11)\cdot 1 = 2(x^2+y-11),
\]
\[
\frac{\partial}{\partial y}\big[(x+y^2-7)^2\big] = 2(x+y^2-7)\cdot 2y = 4y(x+y^2-7),
\]
\[
\frac{\partial f_2}{\partial y} = 2(x^2+y-11) + 4y(x+y^2-7).
\]
Se verificó de forma independiente con WolframAlpha que ambas derivadas
coinciden con esta forma cerrada.

\subsection{Búsqueda de puntos críticos con PyTorch}

Los mínimos observados sobre la malla discreta (\texttt{torch.argmin}) son
solo aproximaciones limitadas por su resolución: para $f_0$ y $f_1$ el
mínimo muestreado se ubicó en $(0,0)$, con valores del orden de
$10^{-13}$--$10^{-14}$, y para $f_2$ en $(3,2)$, con un valor del orden de
$10^{-13}$.

Para refinar esas aproximaciones se calculó el gradiente y la Hessiana
$2\times2$ con \texttt{torch.autograd.grad} y se aplicó el método de
Newton
\[
H(p)\,\Delta p = \nabla f(p), \qquad p_{\text{nuevo}} = p-\Delta p,
\]
partiendo de puntos iniciales tomados de las curvas de nivel. Cada punto
estacionario se clasificó según el signo del determinante $D$ de la
Hessiana: $D<0$ indica un punto silla; $D>0$ con $f_{xx}>0$ indica un
mínimo local; $D>0$ con $f_{xx}<0$ indica un máximo local.

\subsubsection{$f_0$}
$f_0$ tiene un único mínimo global en $(0,0)$, con valor $0$ y
determinante de la Hessiana igual a $4$. Al ser una suma de cuadrados, no
existen otros puntos críticos.

\subsubsection{$f_1$ (Ackley)}
El mínimo global se ubica en $(0,0)$ con valor $0$; ese punto no es
diferenciable en el sentido clásico por el término con raíz cuadrada, por
lo que no se le aplicó la prueba de la Hessiana. Alrededor del origen se
identificaron, mediante Newton, 8 mínimos locales representativos: 4 sobre
los ejes, con valor $\approx 2.5799$ y determinante $\approx 2613.71$; y 4
en las diagonales, con valor $\approx 3.5745$ y determinante $\approx
2623.10$. También se identificaron 4 puntos silla representativos, todos
con valor $\approx 3.2525$ y determinante $\approx -601.48$.

Además, recorriendo cada punto entero no nulo de $[-10,10]^2$ como inicio
de Newton, se identificaron dentro del dominio $440$ mínimos locales
suaves distintos del origen (verificado con un \texttt{assert} en el
notebook). Este conteo es un resultado numérico, no una prueba analítica de
exhaustividad.

\subsubsection{$f_2$ (Himmelblau)}
Se encontraron 4 mínimos globales, todos con valor $0$: $(3, 2)$,
$(-2.805118,\ 3.131313)$, $(-3.779310,\ -3.283186)$ y $(3.584428,\
-1.848127)$. También se identificaron 4 puntos silla, con valores
$13.311926$, $67.719150$, $104.015163$ y $178.337239$ respectivamente.

% TODO: Insertar las figuras de plot_critical_points (mínimos y puntos
% silla marcados sobre las curvas de nivel) exportadas a
% results/figures/.
